% M08, 11.09.2026: F1-Anschluss; Rückweg und Projektion getrennt.
% GithubInboundHarvest/PipelineFullRunner: Sammlung deterministisch,
% fachliche Aufbereitung bedingt modellgestützt; kein automatischer Core-Write.
\section{Projektionen und kontrollierte Außenkopplung}
\label{sec:projektionen-kontrollierte-aussenkopplung}

Der Core wird über operative Sichten nutzbar, ohne seine Rolle als
maßgeblicher fachlicher Bestand abzugeben (Anforderung A8).
Im Besuchsfall F1 erscheinen die Arbeitspakete PBI-046 und PBI-047
als GitHub-Issues \#44 und \#45. Die Projektion macht die geplante
Arbeit außerhalb des Systems zugänglich; sie implementiert die
Besuchsfunktion nicht. Abbildung~\ref{fig:projektion-rueckfuehrung}
zeigt den Hinweg durchgezogen und den kontrollierten Rückweg
gestrichelt.

\begin{figure}[htbp]
  \centering
  \includegraphics[width=\textwidth]{figures/05/abbildung5.7.pdf}
  \caption{Projektions- und Rückführungsmodell der kontrollierten
    Außenkopplung (H: menschlicher Entscheidungspunkt; eigene
    Darstellung).}
  \label{fig:projektion-rueckfuehrung}
\end{figure}

\emph{Projektion.} Ein Erzeugungsvertrag legt fest, wie Titel,
Arbeitsbeschreibung, Akzeptanzkriterien und fachliche Bezüge aus dem
Core als Issue dargestellt werden. Die Inhaltsbildung erfolgt
deterministisch. Die Zuordnung ist auf höchstens ein Issue je
Arbeitspaket und höchstens ein Arbeitspaket je Issue begrenzt;
nicht jedes Arbeitspaket muss bereits projiziert sein. Vor einer
Veröffentlichung wird der Außenbestand als Schnappschuss erhoben.
Die vorgeschlagenen Operationen werden auf diesen Bestand bezogen
und menschlich freigegeben. Die Sicherungen gegen fremde
Schnappschüsse und das Überschreiben externer Änderungen
konkretisiert Abschnitt~\ref{sec:realisierung-persistenz-aussenkopplung};
ihren Gestaltungsanlass erläutert
Abschnitt~\ref{sec:projektionen-aussenkopplung}.

\emph{Rückweg.} Der Abruf und die Auswertung externer Rückmeldungen
werden über den Direktzugang oder nach Werkzeugzustimmung über den
Steward gestartet. Im regulären GitHub-Gesamtlauf werden Issues und
Kommentare neu abgerufen. Die in der Abbildung als \emph{Ernte}
bezeichnete Sammlung erkennt zunächst verarbeitbare Änderungen.
Modellgestützte Komponenten deuten die betreffenden Inhalte und
formulieren Vorschläge mit Herkunftsangaben, das dort bezeichnete
\emph{Destillat}. Diese werden als strukturiertes Änderungspaket
(\emph{Delta}) an den Fortschreibungsablauf übergeben. Ohne
verarbeitbares Änderungspaket endet der Lauf davor. Einordnung,
menschliche Autorisierung und kontrollierte Anwendung bestimmen
erst anschließend die fachliche Zustandsänderung.

Eine spätere, getrennt dokumentierte Kommentarserie greift die
Besuchsübersicht wieder auf: Ein Kommentar fordert eine Anzeige der
nächsten 14 Tage. Laut dem Fallbeleg wurde daraus nach Freigabe eine
zusätzliche Anforderung und eine Aktualisierung des zugehörigen
Arbeitspakets und Issues. Abschnitt~\ref{sec:eval-fortschreibung}
bewertet den Rückweg als F5b und weist seine teilweise narrativen
Belege und weiteren Teilfälle getrennt aus. Der Anschluss zeigt,
wie dasselbe Thema weiterbearbeitet wird; er erweitert F1 nicht
rückwirkend zu einem einzigen durchgehenden Lauf.

\emph{Dokument- und Architektursichten.} Auch sie werden nach außen
veröffentlicht, über dieselbe Freigabeinstanz wie Issues. Bei
hybriden Dokumenten, etwa C4- und Story-Map-Sichten, bleiben
redaktionelle Bereiche von den deterministisch gepflegten
Abschnitten getrennt. Externe Dokumentänderungen werden nicht in
den Core zurückgeführt; Diskussionen und Änderungswünsche nutzen
die vorgesehenen Issue-Kanäle. Fremdänderungen werden bei einem
erneuten Veröffentlichen benannt und bewusst überschrieben.
Architekturinformationen erhalten dabei konkrete Ausgaben:
Begrenzungsbeziehungen liefern technische Rahmenbedingungen in
den betroffenen Issues; als Entscheidungen eingeordnete
Architekturinhalte werden zusätzlich als Architecture Decision
Records dargestellt. Einen dokumentierten Ergebnisfall ordnet
Anhang~\ref{anh:evaluation} zu.

Dieser fachliche Kreislauf verbindet mehrere Läufe. Eine
Prüf-Reparatur-Schleife bearbeitet dagegen ein Zwischenergebnis
(Abschnitt~\ref{sec:saeule-semantische-transformation}); eine
Wiederaufnahme setzt einen pausierten Lauf fort
(Abschnitt~\ref{sec:durable-ausfuehrung-agentische-bedienung}).
Kapitel~\ref{chap:evaluation} untersucht ausgewählte Übergaben und
Kontrollfälle, keine vollständige Erkennung beliebiger externer
Änderungen. Das Ausführungs- und Bedienmodell behandelt der
folgende Abschnitt.
